% Ejercicio 12, restricción — issue #10. Mismo contenido que restriccion.md.
% Fragmento sin preámbulo: lo incluye latex/main.tex, que ya pone el título.
% \sloppy dentro de un grupo: los nombres largos en \texttt no se salen del margen,
% y el ajuste no afecta a lo que sigue.
\begingroup\sloppy

\begin{quote}
\itshape
Las operaciones deben poder ejecutarse en paralelo sobre diferentes cuentas (considera
concurrencia y estado mutable del Visitor). Además, el sistema debe garantizar que las
transacciones sobre una misma cuenta sean atómicas y consistentes (¿cómo lo resuelves con el
patrón elegido?).
\end{quote}

Se construye sobre el código de los incisos~b y~c. Todo el cambio está en
\texttt{ProcesadorTransacciones} (el cliente del patrón) y en un candado por cuenta:
\textbf{ningún visitante cambió}.

\subsubsection{Respuesta corta}

\begin{table}[H]
\centering
\small
\begin{tabularx}{\textwidth}{>{\raggedright\arraybackslash}p{5cm} >{\raggedright\arraybackslash}X}
\toprule
\textbf{Pregunta} & \textbf{Respuesta} \\
\midrule
¿Cómo corren en paralelo las operaciones sobre cuentas distintas?
  & \texttt{EjecutarEnParalelo} reparte las transacciones con \texttt{Parallel.For}. Cada
    transacción solo bloquea sus propias cuentas, así que las de cuentas distintas no se
    esperan entre sí. \\
\addlinespace
¿Qué pasa con el estado mutable del visitante?
  & \texttt{EjecutarTransaccion} y \texttt{RevertirTransaccion} no tienen estado mutable y
    se pueden compartir. \texttt{AuditarTransaccion}, \texttt{CalcularIntereses} y
    \texttt{ResumirPorGrupo} acumulan resultados: se usa \textbf{una instancia por tarea} y
    los resultados se combinan al final. \\
\addlinespace
¿Cómo se garantiza que las transacciones sobre una misma cuenta sean atómicas y consistentes?
  & \textbf{Un candado por cuenta.} Lo toma el cliente del patrón, que envuelve las visitas,
    \textbf{en un orden fijo} (por número de cuenta) y no lo suelta hasta terminar, incluida
    la reversión si algo falla. \\
\addlinespace
¿Cómo se resuelve con el patrón elegido?
  & Visitor separa las reglas de negocio (los visitantes) de quién recorre las cuentas (el
    cliente). La concurrencia es responsabilidad de ese cliente: los visitantes siguen sin
    saber nada de hilos. \\
\bottomrule
\end{tabularx}
\caption{Respuesta corta a la restricción.}
\end{table}

\subsubsection{Paralelismo sobre cuentas distintas}

\begin{lstlisting}[basicstyle=\ttfamily\footnotesize, columns=fullflexible, keepspaces=true]
public IReadOnlyList<ResultadoTransaccion> EjecutarEnParalelo(IReadOnlyList<Transaccion> transacciones)
{
    ResultadoTransaccion[] resultados = new ResultadoTransaccion[transacciones.Count];
    Parallel.For(0, transacciones.Count, i => resultados[i] = Ejecutar(transacciones[i]));
    return resultados;
}
\end{lstlisting}

\begin{itemize}
  \item \textbf{Cada \texttt{Ejecutar} toma solo los candados de sus cuentas.} Dos
        transacciones sobre cuentas distintas no comparten ningún candado y corren a la vez.
        Dos que comparten una cuenta se ejecutan de una en una.
  \item \textbf{Los resultados vuelven en el orden de la lista,} aunque los hilos terminen en
        cualquier orden: cada tarea escribe en su posición del arreglo.
  \item \textbf{Las operaciones sobre un grupo (inciso~c) también corren en paralelo.}
        \texttt{EjecutarEnCadaCuenta} crea una transacción por cuenta del grupo y las ejecuta
        con \texttt{EjecutarEnParalelo}. Como cada una toca una sola cuenta, ninguna espera a
        otra.
\end{itemize}

\subsubsection{El estado mutable de los visitantes}

GoF describe al visitante concreto como el que guarda el estado local del algoritmo, que a
menudo acumula resultados durante el recorrido (Gamma et al., 1994, p.~369). Ese estado es el
que no se puede compartir entre hilos:

\begin{table}[H]
\centering
\small
\begin{tabularx}{\textwidth}{>{\raggedright\arraybackslash}p{3.6cm}
                             >{\raggedright\arraybackslash}X
                             >{\raggedright\arraybackslash}p{2.8cm}
                             >{\raggedright\arraybackslash}X}
\toprule
\textbf{Visitante} & \textbf{Estado mutable} & \textbf{¿Compartible?} & \textbf{Solución} \\
\midrule
\texttt{EjecutarTransaccion} & ninguno: solo guarda la transacción, que no cambia & sí & se
  crea uno por transacción de todos modos \\
\texttt{RevertirTransaccion} & ninguno & sí & ídem \\
\texttt{AuditarTransaccion} & la lista de registros & \textbf{no} & una instancia por tarea
  (\texttt{AuditarEnParalelo}) \\
\texttt{CalcularIntereses} & el diccionario de intereses & \textbf{no} & una instancia por
  recorrido \\
\texttt{ResumirPorGrupo} & los acumuladores por grupo & \textbf{no} & una instancia por
  recorrido \\
\texttt{RecolectarCuentas} & la lista de cuentas & \textbf{no} & \texttt{RecolectarCuentas.De}
  crea uno en cada llamada \\
\bottomrule
\end{tabularx}
\caption{Estado mutable de cada visitante y cómo se evita la condición de carrera.}
\end{table}

\textbf{Qué pasaría si se compartiera un \texttt{AuditarTransaccion} entre hilos:} dos hilos
llamando a \texttt{List.Add} a la vez pueden perder registros, sobrescribirse o lanzar una
excepción. Lo mismo pasa con \texttt{+=} en un acumulador: leer, sumar y escribir son tres
pasos, y dos hilos pueden leer el mismo valor viejo.

\textbf{Solución elegida: una instancia por tarea, y combinar al final.}

\begin{lstlisting}[basicstyle=\ttfamily\footnotesize, columns=fullflexible, keepspaces=true]
public IReadOnlyList<RegistroAuditoria> AuditarEnParalelo(IReadOnlyList<Transaccion> transacciones)
{
    IReadOnlyList<RegistroAuditoria>[] porTarea = new IReadOnlyList<RegistroAuditoria>[transacciones.Count];
    Parallel.For(0, transacciones.Count, i => porTarea[i] = Auditar(transacciones[i]));  // cada Auditar crea su visitante
    return porTarea.SelectMany(registros => registros).ToList();                         // se combinan al final, en orden
}
\end{lstlisting}

Se descartó usar colecciones concurrentes (\texttt{ConcurrentBag}) dentro del visitante por
dos razones:
\begin{itemize}
  \item \textbf{Perderían el orden de los registros,} y la auditoría lo necesita.
  \item \textbf{No resolverían lo importante.} El estado compartido que de verdad importa no
        está en los visitantes, sino en las \textbf{cuentas} (saldo, contadores e historial),
        y eso lo protege el candado del punto siguiente.
\end{itemize}

\subsubsection{Atomicidad y consistencia en la misma cuenta}

\paragraph{Un candado por cuenta.}
Cada \texttt{Cuenta} tiene un objeto privado, \texttt{Candado}, y
\texttt{ProcesadorTransacciones} lo toma con \texttt{Monitor} (lo que usa \texttt{lock} en
C\#) antes de pasar el visitante.
\begin{itemize}
  \item \textbf{Es un objeto propio y no la cuenta misma,} para que ningún código de fuera
        pueda bloquearlo por accidente.
  \item \textbf{Todo lo que cambia la cuenta pasa con el candado tomado:} aplicar el
        movimiento, cobrar el cargo, sumar el contador y anotar el asiento. Ningún otro hilo
        ve la cuenta a medio cambiar.
\end{itemize}

\paragraph{Dónde se toma: en el cliente que envuelve las visitas, no en \texttt{Aceptar}.}
Esta es la respuesta a «¿cómo lo resuelves con el patrón elegido?». Visitor ya separa dos
responsabilidades: los visitantes tienen las reglas de cada tipo de cuenta, y el cliente decide
qué cuentas visitar y con qué visitante. La atomicidad es una propiedad de la
\textbf{transacción entera}, que puede tocar varias cuentas, así que el candado lo toma quien
conoce todas esas cuentas: el cliente.

\begin{lstlisting}[basicstyle=\ttfamily\footnotesize, columns=fullflexible, keepspaces=true]
public ResultadoTransaccion Ejecutar(Transaccion transaccion) =>
    ConCandados(transaccion, () => EjecutarConCandados(transaccion));   // candados alrededor de todas las visitas
\end{lstlisting}

Tomar el candado dentro de \texttt{Aceptar} de cada cuenta no alcanzaría, por tres razones:
\begin{enumerate}
  \item \textbf{Una transferencia necesita las dos cuentas a la vez.} Con el candado en
        \texttt{Aceptar}, la cuenta de origen se soltaría después del débito y antes del
        crédito, y otro hilo podría ver el débito hecho y el crédito no. Si el crédito se
        rechaza, la reversión del débito correría sin protección.
  \item \textbf{No se podría fijar el orden de los candados.} Cada \texttt{Aceptar} tomaría el
        suyo en el orden de los movimientos, que es justo lo que provoca interbloqueos (punto
        siguiente).
  \item \textbf{Mezclaría responsabilidades.} Las cuentas (elementos) y los visitantes
        (reglas) no tienen por qué saber de hilos. Con el candado en el cliente,
        \texttt{EjecutarTransaccion} sigue suponiendo que tiene la cuenta para él solo, y eso
        es cierto.
\end{enumerate}

\paragraph{Orden fijo: sin interbloqueo.}

\begin{lstlisting}[basicstyle=\ttfamily\footnotesize, columns=fullflexible, keepspaces=true]
Cuenta[] cuentas = [.. transaccion.Cuentas.OrderBy(c => c.Numero, StringComparer.Ordinal)];
// se toman en ese orden, se ejecuta la acción y se sueltan en orden inverso
\end{lstlisting}

\textbf{Qué pasaría sin el orden fijo:} dos transferencias cruzadas, A$\rightarrow$B y
B$\rightarrow$A, al mismo tiempo. La primera toma A y espera B; la segunda toma B y espera A.
Ninguna avanza nunca.

\textbf{Con el orden fijo:} las dos piden primero el candado de la cuenta de número menor. La
que lo consigue toma también el otro, termina y suelta los dos; la otra espera y después
sigue. Así se elimina la espera circular, la condición que hace falta para que haya
interbloqueo. Por eso el número de cuenta es único.

\paragraph{Si falla a mitad: todo o nada con \texttt{RevertirTransaccion}.}
Si el débito ya se aplicó y el crédito se rechaza, el procesador deshace el débito con el mismo
\texttt{RevertirTransaccion} del inciso~b, \textbf{todavía con los candados tomados}. Nadie
llega a ver el estado intermedio: desde fuera, la transacción se aplicó entera o no se aplicó.

Por indicación de la revisión del inciso~b, si deshacer también fallara, el procesador ya no
tapa el error original: lanza una \texttt{AggregateException} con los dos errores.

\paragraph{Lecturas consistentes.}
\begin{itemize}
  \item \textbf{\texttt{Auditar} y \texttt{Revertir} también toman los candados de sus
        cuentas,} así que no leen ni deshacen una cuenta a medio cambiar.
  \item \textbf{El historial no se puede recorrer mientras otra transacción escribe en la
        cuenta.} El iterador del inciso~c falla si la lista cambia durante el recorrido, en
        vez de devolver datos a medias. Para leerlo con transacciones en curso hay que hacerlo
        con el candado de la cuenta tomado, como hace la auditoría.
\end{itemize}

\subsubsection{Límites y alternativas}

\begin{itemize}
  \item \textbf{\texttt{Monitor} es síncrono.} Los visitantes no esperan nada (solo hacen
        cuentas), así que alcanza. Si una regla necesitara \texttt{await}, por ejemplo para
        consultar un servicio externo, haría falta \texttt{SemaphoreSlim(1, 1)} con
        \texttt{WaitAsync}, porque un \texttt{Monitor} no se puede mantener a través de un
        \texttt{await}.
  \item \textbf{El candado es por cuenta, no global.} Una cuenta muy activa hace esperar solo a
        sus propias transacciones. Un candado global sería correcto, pero quitaría todo el
        paralelismo.
\end{itemize}

\subsubsection{Demostración}

La sección 7 de la simulación lanza 1000 transacciones a la vez en cada caso. La salida no
depende del orden en que corran los hilos:

\begin{lstlisting}[language={}, numbers=none, basicstyle=\ttfamily\scriptsize, columns=fullflexible, keepspaces=true, breaklines=true]
--- 7. Restricción: transacciones en paralelo ---
  1000 depósitos y retiros de 10.00 a la vez sobre COR-101: 1000 aplicados, saldo 10000.00 (esperado 10000.00), 1000 asientos
  1000 transferencias cruzadas entre COR-101 y COR-102 a la vez: 1000 aplicadas, sin interbloqueo; suma de las dos cuentas 20000.00 (se conserva)
  Auditoría en paralelo de esas transferencias, un AuditarTransaccion por tarea: 2000 registros (2 por transferencia), 2000 conformes
\end{lstlisting}

\texttt{dotnet test} ejecuta \textbf{62 pruebas}, todas correctas: las 57 de los incisos~b
y~c y 5 de concurrencia. Las ejecuté 5 veces seguidas sin ningún fallo.

\begin{table}[H]
\centering
\footnotesize
\begin{tabularx}{\textwidth}{>{\raggedright\arraybackslash}p{7cm} >{\raggedright\arraybackslash}X}
\toprule
\textbf{Prueba} & \textbf{Qué comprueba} \\
\midrule
\texttt{DepositosYRetirosEnParalelo\-SobreLaMismaCuenta\_DanElSaldoEsperado} & 1000 depósitos
  y retiros a la vez sobre \textbf{la misma} cuenta: el saldo final es el esperado y no se
  pierde ningún asiento. \\
\addlinespace
\texttt{TransferenciasCruzadasEnParalelo\_\-TerminanSinInterbloqueoYConservanElDinero} & 1000
  transferencias A$\rightarrow$B y B$\rightarrow$A a la vez terminan antes de 30\,s (sin
  interbloqueo) y la suma de las dos cuentas no cambia. \\
\addlinespace
\texttt{TransferenciasRechazadasEnParalelo\_\-NoDejanRastroEnLaCuentaDeOrigen} & 500
  transferencias que fallan a mitad, mezcladas con 500 depósitos: el origen solo refleja los
  depósitos (todo o nada bajo concurrencia). \\
\addlinespace
\texttt{ReversionesEnParaleloSobreLaMismaCuenta\_\-DejanElEstadoOriginal} & 1000 retiros y sus
  1000 reversiones en paralelo devuelven el saldo, las comisiones y las autorizaciones
  originales. \\
\addlinespace
\texttt{AuditoriasEnParalelo\_ConUnVisitante\-PorTarea\_NoPierdenRegistros} & 1000 auditorías
  en paralelo dan los 1000 registros, en orden. \\
\bottomrule
\end{tabularx}
\caption{Pruebas de concurrencia.}
\end{table}

\paragraph{Referencias.}
Gamma, E., Helm, R., Johnson, R., y Vlissides, J. (1994). \emph{Design patterns: Elements of
reusable object-oriented software}. Addison-Wesley. Los números de página corresponden a la
copia del libro en la bibliografía de la asignatura (\emph{Documentacion Patrones}), igual que
en los incisos anteriores.

\endgroup
